\begin{description}
\item[(a)] Assuming the positions of the two stars are: Vega:
  $(\alpha_1, \delta_1)$; Altair: $(\alpha_2, \delta_2)$
  \[ (\alpha_1, \delta_1) = (279.23538^\circ,\ 38.78547^\circ)
     \quad\text{and}\quad
     (\alpha_2, \delta_2) = (297.69875^\circ,\ 8.87036^\circ) \]
  \hfill(2 points)

  Cosine formula of the spherical geometry is given by:
  \[ \cos[a] = \cos[b]\cos[c] + \sin[b]\sin[c]\cos[A] \]
  \hfill(0.5 point)

  In our triangle, $b = (90 - \delta_1)$, $c = (90 - \delta_2)$ and
  $A = (\alpha_2 - \alpha_1)$. Then the angular distance $\beta$ between
  Vega and Altair is:
  \[ \cos[\beta] = \cos[90^\circ - \delta_1]\cos[90^\circ - \delta_2]
     + \sin[90^\circ - \delta_1]\sin[90^\circ - \delta_2]
       \cos[\alpha_2 - \alpha_1]; \]
  \hfill(3 points)
  \[ \cos[\beta] = \cos[51.21453^\circ]\cos[81.12964^\circ]
     + \sin[51.21453^\circ]\sin[81.12964^\circ]\cos[18.46337^\circ]; \]
  \hfill(1.5 points)
  \[ \beta = 34.19582^\circ \]
  \hfill(2 points)

\item[(b)] The distance of Vega is:
  $r_1 = 1/(130.23\times 10^{-3}) = 7.6787\,pc$

  The distance of Altair is:
  $r_2 = 1/(194.95\times 10^{-3}) = 5.1295\,pc$ \hfill(2 points)

  Using the Law of Cosines:
  $d^2 = r_1^{\,2} + r_2^{\,2} - 2r_1 r_2\cos[\beta]$ \hfill(2 point)

  we can obtain the distance: $d = 4.4855\,pc$ \hfill(2 points)

\item[(c)] Given the proper motions and radial velocities of the two stars,
  the directions of their movements on the celestial sphere can be
  estimated.

  For Vega:
  \begin{align*}
  \mu_{\alpha 1}\cos\delta_1 &= 200.94;\ \mu_{\delta 1} = 286.23\\
  \theta_1 &= \arctan\left[\frac{\mu_{\alpha 1}\cos\delta_1}
    {\mu_{\delta 1}}\right]\times\frac{180^\circ}{\pi} = 35.0697^\circ
  \end{align*}
  \hfill(2 points)

  For Altair:
  \begin{align*}
  \mu_{\alpha 2}\cos\delta_2 &= 536.23;\ \mu_{\delta 2} = 385.29\\
  \theta_2 &= \arctan\left[\frac{\mu_{\alpha 2}\cos\delta_2}
    {\mu_{\delta 2}}\right]\times\frac{180^\circ}{\pi} = 54.302^\circ
  \end{align*}
  \hfill(1 points)

\item[(d)] As the position angles of the proper motion of the two stars are
  different, the stars' paths will intersect. As the paths are circular on
  the celestial sphere, they will intersect in exactly \textbf{two} points.
  \hfill(2 points)

\item[(e)] Let the closer point of intersection be I with the coordinates
  $(\alpha_3, \delta_3)$.

  In triangle PVA,
  \[ PAV = \arcsin\left[\frac{\sin VPA\,\sin PV}{\sin VA}\right]
     = \arcsin\left[\frac{\sin(\alpha_1 - \alpha_2)
       \sin(90^\circ - \delta_1)}{\sin\beta}\right] = 26.055^\circ \]
  \[ PVA = \arcsin\left[\frac{\sin VPA\,\sin PA}{\sin VA}\right]
     = \arcsin\left[\frac{\sin(\alpha_1 - \alpha_2)
       \sin(90^\circ - \delta_2)}{\sin\beta}\right] = 146.196^\circ \]
  \hfill(5 points)

  Note that angle PVA is more than $90^\circ$, which will be evident from
  the diagram. \hfill(1 point)

  The motion of the stars on the celestial sphere will always be along some
  great circle. So VAI is a spherical triangle. Using the four parts formula
  for VAI
  \[ \cot VI = \frac{\cos VA\cos IVA + \sin IVA\cot VAI}{\sin VA} \]
  \[ VI = 90^\circ - \arctan\left[\frac{\cos 34.19582^\circ
     \cos(68.900^\circ) + \sin(68.900^\circ)\cot(99.643^\circ)}
     {\sin 34.19582^\circ}\right] = 76.085^\circ \]
  \hfill(5 points)

  by sine rule,
  \[ AI = \arcsin\left[\frac{\sin VI\sin IVA}{\sin VAI}\right]
     = \arcsin\left[\frac{\sin 76.085^\circ\sin 68.900^\circ}
       {\sin 99.643^\circ}\right] = 66.715^\circ \]
  \hfill(2 points counted in next part)

  Now we use triangle PVI: by cosine rule,
  \begin{align*}
  \cos[PI] &= \cos[PV]\cos[VI] + \sin[PV]\sin[VI]\cos[PVI];\\
  \cos[90^\circ - \delta_3] &= \cos[90^\circ - \delta_1]\cos[VI]
    + \sin[90^\circ - \delta_1]\sin[VI]\cos[180^\circ - \theta_1];\\
  \sin[\delta_3] &= \cos[51.215^\circ]\cos[76.085^\circ]
    + \sin[51.215^\circ]\sin[76.085^\circ]\cos[144.930^\circ];
  \end{align*}
  Thus, $\delta_3 = -27.945^\circ$ \hfill(4 points)

  Applying sine rule,
  \[ VPI = \arcsin\left[\frac{\sin PVI\sin VI}{\sin PI}\right]
     = \arcsin\left[\frac{\sin 144.930^\circ\sin 76.085^\circ}
       {\sin 117.945^\circ}\right] = 39.148^\circ \]
  Hence,
  $\alpha_3 = \alpha_1 - VPI = 279.235^\circ - 39.148^\circ = 240.087^\circ
   \approx 16^{\mathrm{h}}0^{\mathrm{m}}21^{\mathrm{s}}$ \hfill(5 points)

  Out of 5 points, 1 point is reserved for realising that the intersection
  point is in the past path.

\item[(f)] The total velocity along the great circle for Vega is
  \[ \mu_1 = \sqrt{(\mu_{\alpha 1}\cos\delta_1)^2 + \mu_{\delta 1}^{\,2}}
     = \sqrt{200.94^2 + 286.23^2} = 349.72\ \text{mas/year} \]
  \hfill(2 points)
  \[ \text{Number of years} = \frac{VI}{\mu_1}
     = \frac{76.085\times 3600}{0.34972} \approx 783200\ \text{years} \]
  It will happen in the year 781200 BCE. \hfill(2 points)

  Similarly, $\mu_2 = 660.30$ and number of years $= 363700$ years.

  It will happen in the year 3617 BCE.
  % sic: source says "3617 BCE"; 363700 years ago would be ~361700 BCE
  \hfill(2 points)

\item[(g)] After 363700 years, Vega would have traversed $35.335^\circ$
  along its path. \hfill(3 points)

  Thus, its separation from Altair will be
  $(76.085^\circ - 35.335^\circ) = 40.750^\circ$ \hfill(2 points)

\item[(h)] In earth centric Cartesian frame, coordinates of Vega (in units
  of parsec) will be:
  \begin{align*}
  x_1 &= r_1\cos\delta_1\cos(360^\circ - \alpha_1) = 0.96062\,\text{pc}\\
  y_1 &= r_1\cos\delta_1\sin(360^\circ - \alpha_1) = 5.90793\,\text{pc}\\
  z_1 &= r_1\sin\delta_1 = 4.80998\,\text{pc}
  \end{align*}
  Similarly, coordinates of Altair will be, $x_2 = 2.35579$,
  $y_2 = 4.48736$\,pc and $z_2 = 0.79097$\,pc. \hfill(6 points)

  Here, we are assuming the North as positive x-direction and pole as
  positive z-direction.

  The direction vectors for velocities in spherical coordinates are
  \[ v_1 = (-13.9, 7.33, 10.45)\quad\text{and}\quad
     v_2 = (-26.9, 19.57, 14.06) \]
  \hfill(4 points)

  The same in Cartesian coordinates will be
  \begin{align*}
  v_{x1} &= v_{r1}\cos\delta_1\cos\alpha_1 - v_{\alpha 1}\sin\alpha_1
    - v_{\delta 1}\sin\delta_1\cos\alpha_1 = 4.45\\
  v_{y1} &= -v_{r1}\cos\delta_1\sin\alpha_1 - v_{\alpha 1}\cos\alpha_1
    + v_{\delta 1}\sin\delta_1\sin\alpha_1 = -18.33\\
  v_{z1} &= v_{r1}\sin\delta_1 + v_{\delta 1}\cos\delta_1 = -0.56
  \end{align*}
  Similarly, Cartesian velocity vector of Altair will be,
  $v_2 = (4.33, -33.85, 9.87)$. \hfill(8 points)

  Lastly, we evaluate the value of the determinant,
  \[ \begin{vmatrix}
     (x_1 - x_2) & (y_1 - y_2) & (z_1 - z_2)\\
     v_{x1} & v_{y1} & v_{z1}\\
     v_{x1} & v_{y1} & v_{z1}
     % sic: source labels the third row with Vega's components again;
     % the numbers below are Altair's
     \end{vmatrix}
   = \begin{vmatrix}
     -1.40 & 1.42 & 4.02\\
     4.45 & -18.33 & -0.56\\
     4.33 & -33.85 & 9.87
     \end{vmatrix} \]
  \[ = 279.82 - 65.81 - 286.48 = -72.47 \]
  As the value of the determinant is non-zero, the direction vectors of
  these two stars do not cross each other. Hence no such point can exist.
  \hfill(4 points)
\end{description}
